\section{度量怎样唯一选出联络}
\label{sec:math-levi-civita-proof}

度量 \(g\) 量出同一点两向量的内积；联络告诉我们怎样跨点搬运向量。
仅要求平行移动保持内积还不能选出唯一联络，因为沿闭路搬运时仍可附加
局部旋转。再要求没有挠率，就把这份自由度固定了。

\begin{definition}[挠率与度量相容]
切丛联络的挠率是
\[
T(X,Y)=\nabla_XY-\nabla_YX-[X,Y].
\]
若 \(T=0\)，称联络无挠。若对任意向量场 \(X,Y,Z\) 有
\[
X[g(Y,Z)]
=g(\nabla_XY,Z)+g(Y,\nabla_XZ),
\]
称联络与度量 \(g\) 相容。
\end{definition}

在坐标基中，\([\partial_i,\partial_j]=0\)，故无挠等价于
\(\Gamma^k{}_{ij}=\Gamma^k{}_{ji}\)。度量相容则等价于
\(\partial_i g_{jk}
=\Gamma^\ell{}_{ij}g_{\ell k}
+\Gamma^\ell{}_{ik}g_{j\ell}\)。
这些是两组可直接求解联络系数的线性方程。

\begin{theorem}[Levi--Civita 联络]
对任意光滑、非退化的对称度量 \(g\)，存在唯一既与 \(g\) 相容又无挠的
切丛联络。它满足 Koszul 公式
\begin{align*}
2g(\nabla_XY,Z)
={}&Xg(Y,Z)+Yg(Z,X)-Zg(X,Y)\\
&+g([X,Y],Z)-g([Y,Z],X)-g([X,Z],Y),
\end{align*}
在坐标基中等价于
\begin{equation}
\Gamma^k{}_{ij}
=\frac12g^{k\ell}
(\partial_i g_{j\ell}+\partial_j g_{i\ell}
-\partial_\ell g_{ij}).
\label{eq:math-levi-civita-coefficients}
\end{equation}
\end{theorem}
\begin{proof}
先证唯一性。对 \(g(Y,Z),g(Z,X),g(X,Y)\) 分别沿
\(X,Y,Z\) 求导，用度量相容展开，再把前两式相加、第三式相减。
无挠性允许把 \(\nabla_YX\) 换成 \(\nabla_XY-[X,Y]\)，
把 \(\nabla_XZ\) 换成 \(\nabla_ZX+[X,Z]\)，
把 \(\nabla_YZ\) 换成 \(\nabla_ZY+[Y,Z]\)。
含 \(\nabla_ZX,\nabla_ZY\) 的项相消，留下
\(2g(\nabla_XY,Z)\) 和上述三个 Lie 括号项。
由于 \(g\) 非退化，它与所有 \(Z\) 的配对唯一决定 \(\nabla_XY\)。

再证存在性。把 Koszul 公式右边视为对 \(Z\) 的线性泛函，
由 \(g\) 非退化定义唯一向量 \(\nabla_XY\)。
将 \(Y\) 换成 \(fY\)，右边对 \(f\) 求导的项化为
\(2X(f)g(Y,Z)\)，于是
\(\nabla_X(fY)=X(f)Y+f\nabla_XY\)；
将 \(X\) 换成 \(fX\)，所有 \(Y(f),Z(f)\) 项两两相消，得到
\(\nabla_{fX}Y=f\nabla_XY\)。
交换 \(X,Y\) 后相减，右边给
\(2g([X,Y],Z)\)，故 \(T=0\)。
对 \(Y,Z\) 交换并相加，Lie 括号项抵消，得到度量相容式。
这验证构造的确是所需联络。
最后在 \(X=\partial_i,Y=\partial_j,Z=\partial_\ell\)
时代入 \([\partial_i,\partial_j]=0\)，再乘 \(g^{k\ell}\)，
得到坐标公式。
\end{proof}

定理中“非退化”不可删：若 \(g(v,\cdot)=0\) 可对非零 \(v\) 成立，
Koszul 公式就不能由与所有 \(Z\) 的配对唯一找回 \(\nabla_XY\)。
而度量签名不影响证明，所以 Lorentz 度量也有 Levi--Civita 联络。

\section{测地线的两种定义}
\label{sec:math-geodesic-variation}

平行移动已由线性常微分方程定义。一条曲线若自己的速度沿自身平行，
就称为仿射参数化的测地线：
\begin{equation}
\nabla_{\dot\gamma}\dot\gamma=0
\quad\Longleftrightarrow\quad
\ddot x^k+\Gamma^k{}_{ij}\dot x^i\dot x^j=0.
\label{eq:math-geodesic-equation}
\end{equation}
在正定度量中，这与能量泛函
\(E[\gamma]=\frac12\int_{t_0}^{t_1}
g_{ij}(x)\dot x^i\dot x^j\,\dd t\) 的固定端点驻值相同。
为看清联系，直接对 Lagrangian
\(L=\frac12g_{ij}\dot x^i\dot x^j\) 使用 Euler--Lagrange 方程：
\[
\frac{\dd}{\dd t}(g_{kj}\dot x^j)
-\frac12\partial_k g_{ij}\dot x^i\dot x^j=0.
\]
展开时间导数并利用 \(\dot x^i\dot x^j\) 的对称性，再乘 \(g^{\ell k}\)，
恰得到式~\eqref{eq:math-geodesic-equation}。
固定端点使分部积分的边界项消失；“驻值”并不保证整条测地线在全局最短。
Lorentz 度量的零测地线长度为零，仍可由仿射方程定义，
因此也不应称作“最短路”。

\begin{exercise}
把 \(g=dr^2+r^2d\theta^2\) 代入
式~\eqref{eq:math-levi-civita-coefficients}，
重新得到上一节的平面极坐标联络系数。
\end{exercise}
